% ------------------------------------------------------------
% LaTeX Template für die DHBW zum Schnellstart!
% Original: https://github.wdf.sap.corp/vtgermany/LaTeX-Template-DHBW
% ------------------------------------------------------------
% ---- Präambel mit Angaben zum Dokument
% LTeX: enabled=false

\documentclass[
	fontsize=12pt,           % Leitlinien sprechen von Schriftgröße 12.
	paper=A4,
	twoside=false,
	listof=totoc,            % Tabellen- und Abbildungsverzeichnis ins Inhaltsverzeichnis
	bibliography=totoc,      % Literaturverzeichnis ins Inhaltsverzeichnis aufnehmen
	titlepage,               % Titlepage-Umgebung anstatt \maketitle
	headsepline,             % horizontale Linie unter Kolumnentitel
	abstract,              % Überschrift einschalten, Abstract muss in {abstract}-Umgebung stehen
]{scrreprt}                  % Verwendung von KOMA-Report
\usepackage[utf8]{inputenc}  % UTF8 Encoding einschalten
\usepackage[ngerman]{babel}  % Neue deutsche Rechtschreibung
% \usepackage[T1]{fontenc}     % Ausgabe von westeuropäischen Zeichen (auch Umlaute)
% \usepackage{fontspec}
\usepackage{microtype}       % Trennung von Wörtern wird besser umgesetzt
\usepackage{lmodern}         % Nicht-gerasterte Schriftarten (bei MikTeX erforderlich)
\usepackage{graphicx}        % Einbinden von Grafiken erlauben
\usepackage{rotating}        % Rotieren von Grafiken ermöglichen
\usepackage{adjustbox}
\usepackage{wrapfig}         % Grafiken fließend im Text
\usepackage{setspace}        % Zeilenabstand \singlespacing, \onehalfspaceing, \doublespacing
\usepackage[
	%showframe,                % Ränder anzeigen lassen
	left=2.7cm, right=2.5cm,
	top=2.5cm,  bottom=2.5cm,
	includeheadfoot,
	a4paper
]{geometry}                      % Seitenlayout einstellen
\usepackage{scrlayer-scrpage}    % Gestaltung von Fuß- und Kopfzeilen
\usepackage[printonlyused]{acronym}             % Abkürzungen, Abkürzungsverzeichnis
\usepackage{titletoc}            % Anpassungen am Inhaltsverzeichnis
\contentsmargin{0.75cm}          % Abstand im Inhaltsverzeichnis zw. Punkt und Seitenzahl
\usepackage[                     % Klickbare Links (enth. auch "nameref", "url" Package)
  hidelinks,                     % Blende die "URL Boxen" aus.
  breaklinks=true                % Breche zu lange URLs am Zeilenende um
]{hyperref}
\usepackage[hypcap=true]{caption}% Anker Anpassung für Referenzen
\urlstyle{same}                  % Aktuelle Schrift auch für URLs
% Anpassung von autoref für Gleichungen (ergänzt runde Klammern) und Algorithm.
% Anstatt "Listing" kann auch z.B. "Code-Ausschnitt" verwendet werden. Dies sollte
% jedoch synchron gehalten werden mit \lstlistingname (siehe weiter unten).
\addto\extrasngerman{%
	\def\equationautorefname~#1\null{Gleichung~(#1)\null}
	\def\lstnumberautorefname{Zeile}
	\def\lstlistingautorefname{Listing}
	\def\algorithmautorefname{Algorithmus}
	% Damit einheitlich "Abschnitt 1.2[.3]" verwendet wird und nicht "Unterabschnitt 1.2.3"
	% \def\subsectionautorefname{Abschnitt}
}

% ---- Abstand verkleinern von der Überschrift 
\renewcommand*{\chapterheadstartvskip}{\vspace*{.5\baselineskip}}

% Hierdurch werden Schusterjungen und Hurenkinder vermieden, d.h. einzelne Wörter
% auf der nächsten Seite oder in einer einzigen Zeile.
% LaTeX kann diese dennoch erzeugen, falls das Layout ansonsten nicht umsetzbar ist.
% Diese Werte sind aber gute Startwerte.
\widowpenalty10000
\clubpenalty10000

% ---- Für das Quellenverzeichnis
\usepackage[
	backend = biber,                % Verweis auf biber
	language = auto,
	style = numeric,                % Nummerierung der Quellen mit Zahlen
	sorting = none,                 % none = Sortierung nach der Erscheinung im Dokument
	sortcites = true,               % Sortiert die Quellen innerhalb eines cite-Befehls
	block = space,                  % Extra Leerzeichen zwischen Blocks
	hyperref = true,                % Links sind klickbar auch in der Quelle
	%backref = true,                % Referenz, auf den Text an die zitierte Stelle
	bibencoding = auto,
	giveninits = true,              % Vornamen werden abgekürzt
	doi=false,                      % DOI nicht anzeigen
	isbn=false,                     % ISBN nicht anzeigen
    alldates=short                  % Datum immer als DD.MM.YYYY anzeigen
]{biblatex}
\addbibresource{library.bib}
\setcounter{biburlnumpenalty}{3000}     % Umbruchgrenze für Zahlen
\setcounter{biburlucpenalty}{6000}      % Umbruchgrenze für Großbuchstaben
\setcounter{biburllcpenalty}{9000}      % Umbruchgrenze für Kleinbuchstaben
\DeclareNameAlias{default}{family-given}  % Nachname vor dem Vornamen
\AtBeginBibliography{\renewcommand{\multinamedelim}{\addslash\space
}\renewcommand{\finalnamedelim}{\multinamedelim}}  % Schrägstrich zwischen den Autorennamen
\DefineBibliographyStrings{german}{
  urlseen = {Einsichtnahme:},                      % Ändern des Titels von "besucht am"
}
\usepackage[babel,german=quotes]{csquotes}         % Deutsche Anführungszeichen + Zitate


% ---- Für Mathevorlage
\usepackage{amsmath}    % Erweiterung vom Mathe-Satz
\usepackage{amssymb}    % Lädt amsfonts und weitere Symbole
\usepackage{MnSymbol}   % Für Symbole, die in amssymb nicht enthalten sind.

% ---- Für chemische Formeln
\usepackage{mhchem}

% ---- Für Quellcodevorlage
\usepackage{scrhack}                    % Hack zur Verw. von listings in KOMA-Script
\usepackage{listings}                   % Darstellung von Quellcode
\usepackage{xcolor}                     % Einfache Verwendung von Farben
% -- Eigene Farben für den Quellcode
\definecolor{JavaLila}{rgb}{0.4,0.1,0.4}
\definecolor{JavaGruen}{rgb}{0.3,0.5,0.4}
\definecolor{JavaBlau}{rgb}{0.0,0.0,1.0}
\definecolor{ABAPKeywordsBlue}{HTML}{6000ff}
\definecolor{ABAPCommentGrey}{HTML}{808080}
\definecolor{ABAPStringGreen}{HTML}{4da619}
\definecolor{PyKeywordsBlue}{HTML}{0000AC}
\definecolor{PyCommentGrey}{HTML}{808080}
\definecolor{PyStringGreen}{HTML}{008080}
% -- Farben für ABAP CDS
\definecolor{CDSString}{HTML}{FF8C00}
\definecolor{CDSKeywords}{HTML}{6000ff}
\definecolor{CDSAnnotation}{HTML}{00BFFF}
\definecolor{CDSComment}{HTML}{808080}
\definecolor{CDSFunc}{HTML}{FF0000}

% -- Default Listing-Styles

\lstset{
	% Das Paket "listings" kann kein UTF-8. Deswegen werden hier 
	% die häufigsten Zeichen definiert (ä,ö,ü,...)
	literate=%
		{á}{{\'a}}1 {é}{{\'e}}1 {í}{{\'i}}1 {ó}{{\'o}}1 {ú}{{\'u}}1
		{Á}{{\'A}}1 {É}{{\'E}}1 {Í}{{\'I}}1 {Ó}{{\'O}}1 {Ú}{{\'U}}1
		{à}{{\`a}}1 {è}{{\`e}}1 {ì}{{\`i}}1 {ò}{{\`o}}1 {ù}{{\`u}}1
		{À}{{\`A}}1 {È}{{\'E}}1 {Ì}{{\`I}}1 {Ò}{{\`O}}1 {Ù}{{\`U}}1
		{ä}{{\"a}}1 {ë}{{\"e}}1 {ï}{{\"i}}1 {ö}{{\"o}}1 {ü}{{\"u}}1
		{Ä}{{\"A}}1 {Ë}{{\"E}}1 {Ï}{{\"I}}1 {Ö}{{\"O}}1 {Ü}{{\"U}}1
		{â}{{\^a}}1 {ê}{{\^e}}1 {î}{{\^i}}1 {ô}{{\^o}}1 {û}{{\^u}}1
		{Â}{{\^A}}1 {Ê}{{\^E}}1 {Î}{{\^I}}1 {Ô}{{\^O}}1 {Û}{{\^U}}1
		{œ}{{\oe}}1 {Œ}{{\OE}}1 {æ}{{\ae}}1 {Æ}{{\AE}}1 {ß}{{\ss}}1
		{ű}{{\H{u}}}1 {Ű}{{\H{U}}}1 {ő}{{\H{o}}}1 {Ő}{{\H{O}}}1
		{ç}{{\c c}}1 {Ç}{{\c C}}1 {ø}{{\o}}1 {å}{{\r a}}1 {Å}{{\r A}}1
		{€}{{\euro}}1 {£}{{\pounds}}1 {«}{{\guillemotleft}}1
		{»}{{\guillemotright}}1 {ñ}{{\~n}}1 {Ñ}{{\~N}}1 {¿}{{?`}}1,
	breaklines=true,        % Breche lange Zeilen um 
	breakatwhitespace=true, % Wenn möglich, bei Leerzeichen umbrechen
	% Symbol für Zeilenumbruch einfügen
	prebreak=\raisebox{0ex}[0ex][0ex]{\ensuremath{\rhookswarrow}},
	postbreak=\raisebox{0ex}[0ex][0ex]{\ensuremath{\rcurvearrowse\space}},
	tabsize=4,                                 % Setze die Breite eines Tabs
	basicstyle=\ttfamily\small,                % Grundsätzlicher Schriftstyle
	columns=fixed,                             % Besseres Schriftbild
	numbers=left,                              % Nummerierung der Zeilen
	%frame=single,                             % Umrandung des Codes
	showstringspaces=false,                    % Keine Leerzeichen hervorheben
	keywordstyle=\color{blue},
	ndkeywordstyle=\bfseries\color{darkgray},
	identifierstyle=\color{black},
	commentstyle=\itshape\color{JavaGruen},   % Kommentare in eigener Farbe
	stringstyle=\color{JavaBlau},             % Strings in eigener Farbe,
	captionpos=b,                             % Bild*unter*schrift
	xleftmargin=5.0ex
}

% ---- Eigener JAVA-Style für den Quellcode
\renewcommand{\ttdefault}{pcr}               % Schriftart, welche auch fett beinhaltet
\lstdefinestyle{EigenerJavaStyle}{
	language=Java,                             % Syntax Highlighting für Java
	%frame=single,                             % Umrandung des Codes
	keywordstyle=\bfseries\color{JavaLila},    % Keywords in eigener Farbe und fett
	commentstyle=\itshape\color{JavaGruen},    % Kommentare in eigener Farbe und italic
	stringstyle=\color{JavaBlau}               % Strings in eigener Farbe
}

% ---- Eigener ABAP-Style für den Quellcode
\renewcommand{\ttdefault}{pcr}
\lstdefinestyle{EigenerABAPStyle}{
	language=[R/3 6.10]ABAP,
	morestring=[b]\|,                          % Für Pipe-Strings
	morestring=[b]\`,                          % für Backtick-Strings
	keywordstyle=\bfseries\color{ABAPKeywordsBlue},
	commentstyle=\itshape\color{ABAPCommentGrey},
	stringstyle=\color{ABAPStringGreen},
	tabsize=2,
	morekeywords={
		types,
		@data,
		as,
		lower,
		start,
		selection,
		order,
		by,
		inner,
		join,
		key,
		end,
		cast
	}
}

% ---- Eigener Python-Style für den Quellcode
\renewcommand{\ttdefault}{pcr}
\lstdefinestyle{EigenerPythonStyle}{
	language=Python,
	columns=flexible,
	keywordstyle=\bfseries\color{PyKeywordsBlue},
	commentstyle=\itshape\color{PyCommentGrey},
	stringstyle=\color{PyStringGreen}
}

%----- ABAP-CDS-View language
\lstdefinelanguage{ABAPCDS}{
	sensitive=false,
	%Keywords
	morekeywords={define,
		view,
		as,
		select,
		from,
		inner,
		join,
		on,
		key,
		case,
		when,
		then,
		else,
		end,
		true,
		false,
		cast,
		where,
		and,
		distinct,
		group,
		by,
		having,
		min,
		sum,
		max,
		count,
		avg
	},
	%Methoden
	morekeywords=[2]{
		div,
		currency\_conversion,
		dats\_days\_between,
		concat\_with\_space,
		dats\_add_days,
		dats\_is\_valid,
		dats\_add\_months,
		unit\_conversion,
		division,
		mod,
		abs,
		floor,
		ceil,
		round,
		concat,
		replace,
		substring,
		left,
		right,
		length
	},
	morecomment=[s][\color{CDSAnnotation}]{@}{:},
	morecomment=[l][\itshape\color{CDSComment}]{//},
	morecomment=[s][\itshape\color{CDSComment}]{/*}{*/},
	morestring=[b][\color{CDSString}]',
	keywordstyle=\bfseries\color{CDSKeywords},
	keywordstyle=[2]\color{CDSFunc}
}

  % Weitere Details sind ausgelagert

\usepackage{algorithm}                  % Für Algorithmen-Umgebung (ähnlich wie lstlistings Umgebung)
\usepackage{algpseudocode}              % Für Pseudocode. Füge "[noend]" hinzu, wenn du kein "endif",
                                        % etc. haben willst.

\makeatletter                           % Sorgt dafür, dass man @ in Namen verwenden kann.
                                        % Ansonsten gibt es in der nächsten Zeile einen Compilefehler.
\renewcommand{\ALG@name}{Algorithmus}   % Umbenennen von "Algorithm" im Header der Listings.
\makeatother                            % Zeichen wieder zurücksetzen
\renewcommand{\lstlistingname}{Listing} % Erlaubt das Umbenennen von "Listing" in anderen Titel.

% ---- Tabellen
\usepackage{booktabs}  % Für schönere Tabellen. Enthält neue Befehle wie \midrule
\usepackage{multirow}  % Mehrzeilige Tabellen
\usepackage{siunitx}   % Für SI Einheiten und das Ausrichten Nachkommastellen
\sisetup{locale=DE, range-phrase={~bis~}, output-decimal-marker={,}} % Damit ein Komma und kein Punkt verwendet wird.
\usepackage{xfrac} % Für siunitx Option "fraction-function=\sfrac"

% ---- Für Definitionsboxen in der Einleitung
\usepackage{amsthm}                     % Liefert die Grundlagen für Theoreme
\usepackage[framemethod=tikz]{mdframed} % Boxen für die Umrandung
% ---- Definition für Highlight Boxen

% ---- Grundsätzliche Definition zum Style
\newtheoremstyle{defi}
  {\topsep}         % Abstand oben
  {\topsep}         % Abstand unten
  {\normalfont}     % Schrift des Bodys
  {0pt}             % Einschub der ersten Zeile
  {\bfseries}       % Darstellung von der Schrift in der Überschrift
  {:}               % Trennzeichen zwischen Überschrift und Body
  {.5em}            % Abstand nach dem Trennzeichen zum Body Text
  {\thmname{#3}}    % Name in eckigen Klammern
\theoremstyle{defi}

% ------ Definition zum Strich vor eines Texts
\newmdtheoremenv[
  hidealllines = true,       % Rahmen komplett ausblenden
  leftline = true,           % Linie links einschalten
  innertopmargin = 0pt,      % Abstand oben
  innerbottommargin = 4pt,   % Abstand unten
  innerrightmargin = 0pt,    % Abstand rechts
  linewidth = 3pt,           % Linienbreite
  linecolor = gray!40,       % Linienfarbe
]{defStrich}{Definition}     % Name der des formats "defStrich"

% ------ Definition zum Eck-Kasten um einen Text
\newmdtheoremenv[
  hidealllines = true,
  innertopmargin = 6pt,
  linecolor = gray!40,
  singleextra={              % Eck-Markierungen für die Definition
    \draw[line width=3pt,gray!50,line cap=rect] (O|-P) -- +(1cm,0pt);
    \draw[line width=3pt,gray!50,line cap=rect] (O|-P) -- +(0pt,-1cm);
    \draw[line width=3pt,gray!50,line cap=rect] (O-|P) -- +(-1cm,0pt);
    \draw[line width=3pt,gray!50,line cap=rect] (O-|P) -- +(0pt,1cm);
  }
]{defEckKasten}{Definition}  % Name der des formats "defEckKasten"  % Weitere Details sind ausgelagert

% ---- Für Todo Notes
\usepackage{todonotes}
\setlength {\marginparwidth }{2cm}      % Abstand für Todo Notizen

\lstdefinelanguage{JavaScript}{
  morekeywords=[1]{break, case, catch, continue, debugger, default, delete, do, else, finally, for, function, if, in, instanceof, new, return, switch, this, throw, try, typeof, var, void, while, with, require, then, const},
  morecomment=[l]{//},
  morecomment=[s]{/*}{*/},
  morestring=[b]{'},
  morestring=[b]{"},
  morestring=[b]{`},
%   alsoletter={\\'},
  stringstyle=\color[rgb]{0.3,0.5,0.4},
  sensitive=true
}

\lstdefinelanguage{Docker}{
	% Keywords as defined in the language grammar
	morekeywords=[1]{%
	  FROM,RUN,COPY,WORKDIR,LABEL,USER,ENTRYPOINT,EXPOSE,CMD},
	% Built-in functions
	morekeywords=[2]{},
	% Pre-declared types
	morekeywords=[3]{},
	% Constants and zero value
	morekeywords=[4]{},
	% Strings : "foo", 'bar', `baz`
	morestring=[b]{"},
	morestring=[b]{'},
	morestring=[b]{`},
	% Comments : /* cpmment */ and // comment
	comment=[l]{\#},
	morecomment=[s]{/*}{*/},
	stringstyle=\color[rgb]{0.3,0.5,0.4},
	% Options
	sensitive=true
}

\usepackage{xurl}
\usepackage{tikz}

% ---- Elektronische Version oder Gedruckte Version?
% ---- Unterschied: Die elektronische Version enthält keinen Platzhalter für die Unterschrift
\usepackage{ifthen}
\newboolean{e-Abgabe}
\setboolean{e-Abgabe}{false}    % false=gedruckte Fassung

% ---- Persönlichen Daten:
\newcommand{\titel}{Erlernen von Hindernisumfahrung mithilfe von Reinforcement Learning}
\newcommand{\titelheader}{RL-Hindernisumfahrung}
\newcommand{\arbeit}{Studienarbeit (T3\_3101)}
\newcommand{\studiengang}{Informatik}
\newcommand{\studienjahr}{2020}
\newcommand{\autor}{Yannik Schiebelhut}
\newcommand{\autorReverse}{Schiebelhut, Yannik}
\newcommand{\verfassungsort}{Karlsruhe}
\newcommand{\matrikelnr}{3354235}
\newcommand{\kurs}{TINF20B1}
\newcommand{\bearbeitungsmonat}{Mai 2023}
\newcommand{\abgabe}{22. Mai 2023}
\newcommand{\bearbeitungszeitraum}{14.10.2022 - 22.05.2023}
\newcommand{\firmaName}{SAP SE}
\newcommand{\firmaStrasse}{Dietmar-Hopp-Allee 16}
\newcommand{\firmaPlz}{69190 Walldorf, Deutschland}
\newcommand{\betreuerFirma}{Alexander Schaefer}
\newcommand{\betreuerDhbw}{Florian Stöckl}

% ---- Metainformation für das PDF Dokument
\hypersetup{
	pdftitle    = {\titel},
	pdfsubject  = {\arbeit},
	pdfauthor   = {\autor},
	%pdfkeywords = {Keywords angeben},
	pdfcreator  = {LaTeX},
	%pdfproducer = {in der Regel pdfTeX}
}

% ---- Definition der Kopf- und Fußzeilen
\clearscrheadfoot                               % Löschen von LaTeX Standard
\automark[section]{chapter}                     % Füllen von section und chapter
\renewcommand*{\chaptermarkformat}{}            % Entfernt die Kapitelnummer
\renewcommand*{\sectionmarkformat}{}            % Entfernt die Sectionnummer
% Angaben [für "plain"]{für "scrheadings"}
% \ihead[]{\titelheader}                          % Kopfzeile links
\ihead{\ifthenelse{\value{chapter}>0}{\chaptername~\thechapter}{}}
\chead[]{}                                      % Kopfzeile mitte
% \ohead[]{\rightmark}                            % Kopfzeile rechts
\ohead{\ifthenelse{\value{chapter}>0}{\headmark}{}}
\ifoot[]{}                                      % Fußzeile links
\cfoot*{\sffamily\pagemark}                     % Fußzeile mitte
\ofoot[]{}                                      % Fußzeile rechts
\KOMAoptions{
   headsepline = 0.2pt,                         % Liniendicke Kopfzeile
   footsepline = false                          % Liniendicke Fußzeile
}

% ---- Hilfreiches
\newcommand{\zB}{z.\,B. }   % "z.B." mit kleinem Leeraum dazwischen (ohne wäre nicht korrekt)
\newcommand{\dash}{d.\,h. }
\newcommand{\Dash}{D.\,h. }

\newcommand{\code}[1]{\texttt{#1}} % Ist einfacher zu schreiben als ständig \texttt und erlaubt
                                   % Änderungen im Nachhinein, wenn man z.B. Inline-Code anders stylen möchte.

% ---- Silbentrennung (falls LaTeX defaults falsch / nicht gewünscht sind)
\hyphenation{HANA}         % anstatt HA-NA
\hyphenation{Graph-Script} % anstatt GraphS-cript

% ---- Beginn des Dokuments
\begin{document}
\setlength{\parindent}{0pt}              % Keine Paragraphen Einrückung.
                                         % Dafür haben wir den Abstand zwischen den Paragraphen.
\setcounter{secnumdepth}{2}              % Nummerierungstiefe fürs Inhaltsverzeichnis
\setcounter{tocdepth}{1}                 % Tiefe des Inhaltsverzeichnisses. Ggf. so anpassen,
                                         % dass das Verzeichnis auf eine Seite passt.
\sffamily                                % Serifenlose Schrift verwenden.

% ---- Vorspann
% ------ Titelseite
\singlespacing
\thispagestyle{empty}
\begin{titlepage}
\enlargethispage{4cm}

\begin{figure}
	\includegraphics[height=2.5cm]{Bilder/Logos/Logo_DHBW.pdf}
	\centering
\end{figure} 
\vspace*{0.1cm}

\begin{center}
	\huge{\textbf{\titel}}\\[1.5cm]
	\Large{\textbf{\arbeit}}\\[0.5cm]
	\normalsize{im Rahmen der Prüfung zum\\[1ex] \textbf{Bachelor of Science (B.Sc.)}}\\[0.5cm]
	\Large{des Studienganges \studiengang}\\[1ex]
	\normalsize{an der Dualen Hochschule Baden-Württemberg Karlsruhe}\\[1cm]
	\normalsize{von}\\[1ex] \Large{\textbf{\autor}} \\[1cm]
\end{center}

\begin{center}
	\vfill
	\begin{tabular}{ll}
		Abgabedatum:                     & \abgabe \\[0.2cm]
		Bearbeitungszeitraum:            & \bearbeitungszeitraum \\[0.2cm]
		Matrikelnummer, Kurs:            & \matrikelnr , \kurs \\[0.2cm]
		% Ausbildungsfirma:                & \firmaName \\
		%                                  & \firmaStrasse \\
		%                                  & \firmaPlz \\[0.2cm]
		% Betreuer der Ausbildungsfirma:   & \betreuerFirma \\[0.2cm]
		Gutachter der Dualen Hochschule: & \betreuerDhbw \\[2cm]
	\end{tabular} 
\end{center}
\end{titlepage}
  % Titelseite
\newcounter{savepage}
\pagenumbering{Roman}                    % Römische Seitenzahlen
\onehalfspacing

% ------ Erklärung, Sperrvermerk, Abstact
\chapter*{Eidesstattliche Erklärung}
Ich versichere hiermit, dass ich meine \arbeit{} mit dem Thema:
\begin{quote}
	\textit{\titel}
\end{quote} 
gemäß § 5 der \enquote{Studien- und Prüfungsordnung DHBW Technik} vom 29. September 2017 selbstständig verfasst und keine anderen als die angegebenen Quellen und Hilfsmittel benutzt habe. Die Arbeit wurde bisher keiner anderen Prüfungsbehörde vorgelegt und auch nicht veröffentlicht.

\vspace{0.25cm}

Ich versichere zudem, dass die eingereichte elektronische Fassung mit der gedruckten Fassung übereinstimmt.

\vspace{1cm}

\verfassungsort, den \today \\[0.5cm]
\ifthenelse{\boolean{e-Abgabe}}
	{\underline{Gez. \autor}}
	{\makebox[6cm]{\hrulefill}}\\ 
\autorReverse

\renewcommand{\abstractname}{Abstract} % Veränderter Name für das Abstract
\begin{abstract}
\begin{addmargin}[1.5cm]{1.5cm}        % Erhöhte Ränder, für Abstract Look
\thispagestyle{plain}                  % Seitenzahl auf der Abstract Seite

\begin{center}
\small\textit{- Deutsch -}             % Angabe der Sprache für das Abstract
\end{center}

\vspace{0.25cm}

Platzhalter


\end{addmargin}
\end{abstract}
\renewcommand{\abstractname}{Abstract} % Veränderter Name für das Abstract
\begin{abstract}
\begin{addmargin}[1.5cm]{1.5cm}        % Erhöhte Ränder, für Abstract Look
\thispagestyle{plain}                  % Seitenzahl auf der Abstract Seite

\begin{center}
\small\textit{- English -}             % Angabe der Sprache für das Abstract
\end{center}

\vspace{0.25cm}

% LTeX: language=en-US

Placeholder

\end{addmargin}
\end{abstract}

% ------ Inhaltsverzeichnis
\singlespacing
\tableofcontents

% ------ Verzeichnisse
\renewcommand*{\chapterpagestyle}{plain}
\pagestyle{plain}
% \chapter*{Formelverzeichnis}
\addcontentsline{toc}{chapter}{Formelverzeichnis} % Hinzufügen zum Inhaltsverzeichnis 

% Definition des neuen Befehls für das Einfügen der Abkürzung der Einheit
\newcommand{\acrounit}[1]{
  \acroextra{\makebox[18mm][l]{\si[per-mode=fraction,fraction-function=\sfrac]{#1}}}
}
\begin{acronym}[dmin] % längstes Kürzel wird verw. für den Abstand zw. Kürzel u. Text

	% Alphabetisch selbst sortieren - nicht verwendete Formeln rausnehmen!
	% Allgemein: \acro{KÜRZEL}[ABKÜRZUNG]{\acrounit{SI-EINHEIT}BESCHREIBUNG}

	\acro{A}[\ensuremath{A}]{\acrounit{mm^2}Fläche}	
	\acro{D}[\ensuremath{D}]{\acrounit{mm}Werkstückdurchmesser}	
	\acro{dmin}[\ensuremath{d\textsubscript{min}}]{\acrounit{mm}kleinster Schaftdurchmesser}	
	\acro{L1}[\ensuremath{L\textsubscript{1}}]{\acrounit{mm}Länge des Werkstückes Nr. 1}	
	\acro{Fwinkel}[]{\acrounit{Grad}Freiwinkel}	
	\acro{Kwinkel}[]{\acrounit{Grad}Keilwinkel}

\end{acronym}

\chapter*{Abkürzungsverzeichnis}
\addcontentsline{toc}{chapter}{Abkürzungsverzeichnis} % Hinzufügen zum Inhaltsverzeichnis 

\begin{acronym}[PPO] % längstes Kürzel wird verw. für den Abstand zw. Kürzel u. Text

	% Alphabetisch selbst sortieren - nicht verwendete Kürzel rausnehmen!
	% \acro{AIR}{Adobe Integrated Runtime}
	% \acro{AJAX}{Asynchronous Javascript and XML}
	% \acro{ANSI}{American National Standards Institute}
	% \acro{api}[API]{Application Programming Interface}
	% \acro{AR}{Augmented Reality}
	% \acro{BAPI}{Business Application Programming Interface}
	% \acro{BIOS}{Basic Input Output System}
	% \acro{bpmn}[BPMN]{Business Process Model and Notation}
	% \acro{ccwf}[ccWF]{Cross-Company Workflow}
	% \acro{CDMA}{Code Division Multiple Access}
	% \acro{csv}[CSV]{Comma-Separated Values}
	% \acro{erp}[ERP]{Enterprise Resource Planning}
	% \acro{hs2}[HS2]{High Speed 2}
	% \acro{HTTP}{Hypertext Transfer Protocol}
	% \acro{HTTPS}{Hypertext Transfer Protocol Secure}
	% \acro{IP}{Internet Protocol}
	% \acro{ISBN}{Internationale Standardbuchnummer}
	% \acrodefplural{ISBN}[ISBNs]{Internationale Standardbuchnummern}
	% \acro{JSON}{JavaScript Object Notation}
	% \acro{kpi}[KPI]{Key Performance Indicator}
	\acro{mdp}[MDP]{Markov Decision Process}
	% \acro{OData}{Open Data Protocol}
	\acro{ppo}[PPO]{Proximal Policy Optimization}
	% \acro{SDK}{Software Development Kit}
	% \acro{SEO}{Search Engine Optimization}
	% \acro{sql}[SQL]{Structured Query Language}
	% \acro{SSH}{Secure Shell}
	% \acro{UEFI}{Unified Extensible Firmware Interface}
	% \acro{URI}{Uniform Resource Identifier}
	% \acro{USB}{Universal Serial Bus}
	% \acro{VLAN}{Virtual Local Area Network}
	% \acro{WYSISWG}{What You See Is What You Get}
	% \acro{xes}[XES]{eXtensible Event Stream}
	% \acro{XSL}{Extensible Stylesheet Language}

\end{acronym}
\listoffigures                          % Erzeugen des Abbildungsverzeichnisses 
% \listoftables                           % Erzeugen des Tabellenverzeichnisses
\renewcommand{\lstlistlistingname}{Quellcodeverzeichnis}
\lstlistoflistings                      % Erzeugen des Listenverzeichnisses
\setcounter{savepage}{\value{page}}


% ---- Inhalt der Arbeit
\cleardoublepage
\pagenumbering{arabic}                  % Arabische Seitenzahlen für den Hauptteil
\setlength{\parskip}{0.5\baselineskip}  % Abstand zwischen Absätzen
\rmfamily
\renewcommand*{\chapterpagestyle}{scrheadings}
\pagestyle{scrheadings}
\onehalfspacing
% \include{content goes here}
\chapter{Einleitung}
\label{sec:einleitung}
Die Robotik ist ein breites Forschungsfeld mit praktisch grenzenlosen Möglichkeiten.
Die Fähigkeit, sich zu bewegen, ist dabei besonders spannend, da sie die Flexibilität der Einsatzmöglichkeiten für Roboter stark erhöht.
Roboter mit Beinen stellen sich besonders heraus.
Mit Inspirationen aus Mensch- und Tierreich bieten diese das Potenzial, sich in jedem denkbaren Terrain fortzubewegen, das auch für Menschen zugänglich ist oder sogar in Gebiete vorzudringen, die uns verwehrt bleiben.
Im Vergleich zu anderen Fortbewegungsarten stellt die stabile Koordination von mehreren Beinen, die jeweils aus mehreren Gelenken bestehen, allerdings eine große Herausforderung für die technische Umsetzung dar.
In der Regel sind für die Programmierung solcher Roboter eine sehr genaue Kenntnis der Maschine und deren Dynamik vonnöten.

In den vergangenen Jahren wird deshalb verstärkt erforscht, wie Roboter sich diese Fähigkeiten selbstständig mittels Reinforcement Learning beibringen können.
In einer vorangegangenen Studienarbeit wurden am Beispiel eines eigens dafür gebauten, spinnenartigen, vierbeinigen Roboters Möglichkeiten erforscht, um diesen mittels selbst gelernter Bewegungsabläufe möglichst effizient und schnell geradlinig nach vorne zu bewegen.
Um den Trainingsprozess zu beschleunigen wurde dabei der Roboter in der Simulationsumgebung \enquote{Unity} nachgebaut und trainiert.

Für das Erfüllen eines praktischen Nutzens ist es jedoch in der Regel nicht ausreichend, wenn sich ein Roboter nur in eine feste Richtung bewegen kann.
Im Rahmen dieser Arbeit wird deshalb untersucht, wie dem Roboter beigebracht werden kann, einem gezielt übergebenen Pfad zu folgen.
Weiterhin soll der Roboter dabei Hindernisse, die sich auf diesem Pfad befinden, automatisch umsteuern und anschließend wieder auf den vorgebenen Pfad zurückkehren. % ggf Absatz hier
Dazu werden zunächst die Arbeitsumgebung und Ergebnisse der vorherigen Arbeit rekonstruiert.
Anschließend wird diskutiert, welche Änderungen am Roboter und dessen Simulationsumgebung vorgenommen werden müssen, um die erweiterten Anforderungen grundsätzlich erfüllen zu können.
Außerdem wird erläutert, wie die Aufgabe in sinnvolle Teilaufgaben gegliedert werden kann.
Zur Bearbeitung dieser Teilaufgaben wird dann ein Konzept erarbeitet, welches im Anschluss in der Simulationsumgebung, unter Zuhilfenahme des ML-Agents-Toolkits, mittels des \acl{ppo}-Algorithmus ein Proof-of-Concept trainiert.
Anhand der Ergebnisse dieses Trainingsprozesses wird evaluiert, ob das entwickelte Konzept einen erfolgversprechenden Lösungsansatz darstellt und welche Verbesserungen vorgenommen werden sollten.

% \begin{itemize}
%     \item Roboter sind ein breites Forschungsfeld mit grenzenlosen Möglichkeiten
%     \item Fähigkeit, sich zu bewegen ist besonders spannend für die Robotik, da es die Einsatzmöglichkeiten für Roboter stark erweitert
%     \item dabei insbesondere mehrbeinige Roboter interessant
%     \item mit Inspiration aus Menschen- und Tierreich sind diese fähig, sich in praktisch jedem denkbaren Terrain fortzubewegen
%     \item im Vergleich zu anderen Fortbewegungsarten stellt stabile Koordination von Beinen jedoch eine große Herausforderung für die technische Umsetzung voraus
%     \item in der Regel sehr genaue Kenntnis des Roboters und dessen Dynamik vonnöten
    
%     \item in den vergangenen Jahren wird deshalb verstärkt erforscht, wie Roboter sich diese Fähigkeit gegebenenfalls selbstständig mittels Reinforcement Learning beibringen können
%     \item in einer vorangegangenen Studienarbeit wurden am Beispiel eines eigens dafür gebauten, vierbeinigen Roboters Möglichkeiten erforscht, um diesen möglichst effizient und schnell geradlinig nach vorne zu bewegen
%     \item dabei wurde der Roboter in der Simulationsumgebung \enquote{Unity} trainiert
%     \item für das Erfüllen eines praktischen Nutzens ist es jedoch in der Regel nicht ausreichend, wenn sich ein Roboter nur in eine Richtung bewegen kann
%     \item im Rahmen dieser Arbeit wird deshalb untersucht, wie dem Roboter beigebracht werden kann, einem gezielt übergebenen Pfad zu folgen
%     \item weiterhin soll der Roboter Hindernisse, die sich auf diesem Pfad befinden, automatisch umsteuern und anschließend wieder auf den vorgegebenen Pfad zurückkehren
    
%     \item dazu werden zunächst die Ergebnisse und Arbeitsumgebung der vorherigen Arbeit rekonstruiert
%     \item anschließend wird diskutiert, welche Änderungen am Roboter und dessen Simulationsumgebung vorgenommen werden müssen, um die erweiterten Anforderungen grundsätzlich erfüllen zu können
%     \item außerdem wird erläutert, wie die Aufgabe in sinnvolle Teilaufgaben gegliedert werden kann
%     \item zur Bearbeitung dieser Teilaufgaben wird dann ein Konzept erarbeitet, welches im Anschluss in der Simulationsumgebung unter Zuhilfenahme des ML-Agents-Toolkits als Proof-of-Concept trainiert wird
%     \item Proximal-Policy-Optimization
% \end{itemize}


% ---- Literaturverzeichnis
\cleardoublepage
\renewcommand*{\chapterpagestyle}{plain}
\pagestyle{plain}
\pagenumbering{Roman}                   % Römische Seitenzahlen
\setcounter{page}{\numexpr\value{savepage}+1}
\printbibliography[title=Literaturverzeichnis]

% ---- Anhang
\appendix
\chapter{Hinweis}
Der vollständige Quellcode sämtlicher Trainingsschritte ist frei verfügbar unter \url{https://github.com/MobMonRob/HindernisumfahrungRLStudien} sowie unter \url{https://github.com/yschiebelhut/spider_bot_rl_training}.
Weiterhin enthält \url{https://github.com/yschiebelhut/spiderbot-training-data} die für Linux kompilierten Trainingsumgebungen, die trainierten Modelle und alle notwendigen Daten, um die Trainingsprozesse mit Tensorboard nachvollziehen zu können.

\paragraph{Wichtig:} Zur Durchführung der Inferenz eines bestimmten Modells sollte das Haupt-Repository möglichst auf dem entsprechenden Commit ausgecheckt sein, auf dem das Training stattgefunden hat.
(Die Numerierungen in den Commits und den Run-IDs geben darüber Aufschluss.)
Dies ist wichtig, da insbesondere die Größe des Observation Space als auch die Observations selbst mit den beim Training verwendeten Daten übereinstimmen müssen, weil es ansonsten zu Fehlern bei der Interpretation des neuronalen Netzes kommt.


\chapter{ML-Agents Training Konfiguration}
\label{anhang:trainer-config}
\begin{figure}[H]
    \lstinputlisting[
        label = code:trainer-config-yaml
    ]{Code/new-trainer-config.yaml}
\end{figure}
%\clearpage
%\pagenumbering{Roman}  % römische Seitenzahlen für Anhang

\newpage
\end{document}
